% ─────────────────────────────────────────────────────────────────────────────
\section{Le facteur \texorpdfstring{$k_{mod}$}{kmod}~:
         l'effet de la dur\'ee de charge}
% ─────────────────────────────────────────────────────────────────────────────

\begin{maitre}
Pourquoi le bois sous une charge de courte dur\'ee (neige)
r\'esiste mieux que sous une charge permanente~?
\end{maitre}

\begin{apprenti}
\`A cause du fluage~? Le bois se d\'eforme progressivement
sous charge longue, ce qui affaiblit sa r\'esistance~?
\end{apprenti}

\begin{maitre}
Exactement. L'Eurocode~5 traduit cela par le facteur $k_{mod}$ du
Tableau~3.1~: plus la charge est de longue dur\'ee, plus $k_{mod}$
est faible, et donc plus la r\'esistance de calcul $f_{m,d}$ est diminu\'ee.
\end{maitre}

La r\'esistance de calcul $f_{m,d}$ est obtenue \`a partir de la
r\'esistance caract\'eristique $f_{m,k}$ (issue des tables EN~338) par~:

\begin{equation}
f_{m,d} = f_{m,k} \cdot \frac{k_{mod}}{\gamma_M}
\label{eq:fmd}
\end{equation}

\noindent avec $\gamma_M = 1{,}3$ pour le bois massif (EC5 \S{}2.4).

\begin{center}
\begin{tabular}{ll}
\toprule
\textbf{Dur\'ee de charge} & \boldmath{$k_{mod}$} (bois massif, CS~1)\\
\midrule
Permanente & 0,60\\
Long terme & 0,70\\
Moyen terme & 0,80\\
Court terme (neige, vent) & 0,90\\
Instantan\'ee & 1,10\\
\bottomrule
\end{tabular}
\end{center}

Exemple num\'erique pour un bois C24 sous neige (court terme, CS~1)~:
\[
f_{m,d} = 24{,}0 \times \frac{0{,}90}{1{,}3} = 16{,}6 \text{ MPa}
\]

\begin{lstlisting}[style=python, caption={Usage de k\_mod dans \fichier{ec5/elu.py}}]
for combi in combinaisons:
    # Lecture de k_mod depuis la table kmod.csv
    # selon (famille du bois, classe de service, durée de charge)
    k_mod   = get_kmod(famille, classe_service, combi.duree_charge)
    gamma_M = get_gamma_m(famille)   # 1.3 pour bois massif

    # Résistance de calcul en flexion (MPa)
    f_m_d = materiau.f_m_k_MPa * k_mod / gamma_M
    # Ex : 24.0 × 0.90 / 1.3 = 16.6 MPa (neige, CS1)
\end{lstlisting}